\section*{Chapitre \ref{ch:b2:polymer-synthesis} --- Polymères~: synthèse, structure et propriétés}

\begin{solution}{exo:b2:polymer-synthesis:1}
\omterm{def:b2:polymer-synthesis:polymer}{Motifs de répétition}~: polyéthène, \ce{-[CH2-CH2]-}~; polypropène, \ce{-[CH2-CH(CH3)]-}~; poly(chlorure de vinyle),
\ce{-[CH2-CHCl]-}~; polystyrène, \ce{-[CH2-CH(C6H5)]-}~; enfin, pour le nylon-6,6, le motif \ce{-[NH(CH2)6NH-CO(CH2)4CO]-}.
\end{solution}

\begin{solution}{exo:b2:polymer-synthesis:2}
PET~: croissance par étapes (polyester, eau libérée). Polystyrène~: croissance en chaîne. Nylon-6,6~:
croissance par étapes. Poly(méthacrylate de méthyle)~: croissance en chaîne (un \omterm{def:b2:polymer-synthesis:polymer}{monomère} à \ce{C=C}).
Acide polylactique à partir de l'hydroxyacide~: croissance par étapes, un \omterm{def:b2:polymer-synthesis:polymer}{monomère} \ce{A-B}.
\end{solution}

\begin{solution}{exo:b2:polymer-synthesis:3}
$\bar M_n = 1000 \times \qty{104.15}{g/mol} \approx \qty{104}{kg/mol}$ (styrène \ce{C8H8},
\qty{104.152}{g/mol}).
% ledger: aw:C, aw:H
\end{solution}

\begin{solution}{exo:b2:polymer-synthesis:4}
\omterm{def:b2:polymer-synthesis:tacticity}{Syndiotactique}. Oui~: une chaîne régulière peut s'empiler en cristallites.
\end{solution}

\begin{solution}{exo:b2:polymer-synthesis:5}
Fractions massiques 0.5 et 0.5. $1/\bar M_n = 0.5/10 + 0.5/100 = 0.055$, $\bar M_n \approx
\qty{18.2}{kg/mol}$~; $\bar M_w = 0.5 \times 10 + 0.5 \times 100 = \qty{55}{kg/mol}$~; $\text{\DH} \approx
3.0$.
\end{solution}

\begin{solution}{exo:b2:polymer-synthesis:6}
$1/(1-0.98) = 50$~; $1/(1-0.995) = 200$.
\end{solution}

\begin{solution}{exo:b2:polymer-synthesis:7}
$r = 1/1.01$~; pour $p = 1$, $\bar X_n = (1+r)/(1-r) = (1.01 + 1)/(1.01 - 1) = 201$.
\end{solution}

\begin{solution}{exo:b2:polymer-synthesis:8}
$R_p \propto [\mathrm I]^{1/2}$~: multipliée par $\sqrt2 \approx 1.41$. $\nu \propto [\mathrm I]^{-1/2}$~:
divisée par $\sqrt 2$, soit environ 0.71 fois sa valeur antérieure.
\end{solution}

\begin{solution}{exo:b2:polymer-synthesis:9}
Le PVC pur a sa transition vitreuse au-dessus de la température ambiante~: c'est un verre rigide. Les
petites molécules de plastifiant se logent entre les chaînes, les écartent et permettent aux segments de
bouger à plus basse température~: $T_g$ descend
au-dessous de la température ambiante, et le matériau est caoutchouteux et souple à l'usage.
\end{solution}

\begin{solution}{exo:b2:polymer-synthesis:10}
$\dfrac{\mathrm d}{\mathrm dx}\bigl(x p^{\,x-1}\bigr) = p^{\,x-1}(1 + x\ln p)$, nulle pour $x = -1/\ln p$
(un maximum~: le signe passe de $+$ à $-$). Pour $p$ voisin de 1, $\ln p \approx -(1-p)$, donc $x \approx
1/(1-p) = \bar X_n$.
\end{solution}

\begin{solution}{exo:b2:polymer-synthesis:11}
Étape 1, \omterm{def:b2:acyl-substitution:transesterification}{transestérification}~: \ce{C6H4(COOCH3)2 + 2HOCH2CH2OH -> C6H4(COOCH2CH2OH)2 + 2CH3OH}, le
méthanol étant distillé. Étape 2~: chaque molécule de diester porte deux extrémités \ce{CH2CH2OH}~; une
extrémité attaque un ester d'une autre molécule et libère une molécule d'éthane-1,2-diol, éliminée par
pompage. Le \omterm{def:b2:polymer-synthesis:polymer}{monomère} porte les deux sortes de groupes dans le bon rapport, comme un \omterm{def:b2:polymer-synthesis:polymer}{monomère} \ce{A-B}~:
l'équilibre $r = 1$ est
intégré, et l'élimination du diol pousse $p$ vers 1.
\end{solution}

\begin{solution}{exo:b2:polymer-synthesis:12}
Alimentation équimolaire~: rapport $= (2.0 + 1)/(1 + 0.5) = 2$~: deux fois plus d'unités du \omterm{def:b2:polymer-synthesis:polymer}{monomère} 1.
Avec $r_1 = r_2 =
0$, rapport $= [\mathrm M_1][\mathrm M_2]/([\mathrm M_2][\mathrm M_1]) = 1$ quelle que soit l'alimentation~:
chaque radical n'additionne que l'autre \omterm{def:b2:polymer-synthesis:polymer}{monomère}, d'où un \omterm{def:b2:polymer-synthesis:copolymer}{copolymère} alterné.
\end{solution}

\begin{solution}{pb:b2:polymer-synthesis:1}
\textbf{1.} \ce{H2N(CH2)6NH2} et \ce{HOOC(CH2)4COOH}.
\textbf{2.} Une molécule d'eau par liaison~:
\[\mathrm{{-}COOH + H_2N{-} \longrightarrow {-}CO{-}NH{-} + H_2O}.\]
\textbf{3.} Le sel contient exactement une diamine par diacide~; deux pesées séparées
n'atteindraient jamais l'exactitude qu'exige l'équation de Carothers.
\textbf{4.} \ce{-[NH(CH2)6NH-CO(CH2)4CO]-}, \ce{C12H22N2O2}, \qty{226.32}{g/mol}.
\textbf{5.} Le \omterm{def:b2:polymer-synthesis:polymer}{motif de répétition} contient deux unités issues de \omterm{def:b2:polymer-synthesis:polymer}{monomères}~: $M_0 = 226.32/2 = \qty{113.16}{g/mol}$.
\textbf{6.} Les carbones de la diamine (6) et ceux du diacide (6).
\textbf{7.} $\bar X_n = 50$, $\bar M_n = 50 \times 113.16 \approx \qty{5.66}{kg/mol}$.
\textbf{8.} $\bar X_n = 100$, $\bar M_n \approx \qty{11.3}{kg/mol}$.
\textbf{9.} Avec 99~\% des groupes ayant réagi, la chaîne moyenne ne compte que 100 unités, trop peu pour
une fibre solide~; chaque pas supplémentaire vers la \omterm{def:b2:continuous-reactors:conversion}{conversion} totale double ou triple la longueur des chaînes.
\textbf{10.} $\bar X_n = 10\,000/113.16 = 88.4$~; $p = 1 - 1/88.4 \approx 0.989$.
\textbf{11.} En éliminant l'eau~: la masse fondue est chauffée bien au-dessus de sa température de fusion,
finalement sous pression réduite, afin que l'équilibre continue de se déplacer vers l'amide.
\textbf{12.} Une impureté monofonctionnelle bloque des extrémités de chaîne, et toute impureté rompt
l'équilibre 1\,:\,1.
\textbf{13.} $r = 1/1.01 \approx 0.990$.
\textbf{14.} $\bar X_n = (1+r)/(1-r) = 201$~; $\bar M_n \approx 201 \times 113.16 \approx
\qty{22.7}{kg/mol}$.
\textbf{15.} $\bar X_n = 1.9901/(1.9901 - 2 \times 0.9901 \times 0.99) \approx 1.9901/0.0297 \approx
67$.
\textbf{16.} Des groupes amine~: quand les groupes acide sont épuisés, chaque chaîne se termine par \ce{NH2}.
\textbf{17.} Il transforme une extrémité amine en un amide qui ne peut plus réagir~: il bloque des chaînes,
ce qui limite $\bar M_n$ et arrête toute croissance ultérieure quand le \omterm{def:b2:polymer-synthesis:polymer}{polymère} est refondu.
\textbf{18.} Pour garder une viscosité à l'état fondu reproductible pour le filage et le moulage, viscosité
qui dériverait sinon à mesure que les chaînes continuent de réagir dans la masse fondue.
\textbf{19.} $\text{\DH} = 1 + p = 1.99$.
\textbf{20.} $\bar M_w = 1.99 \times 11.3 \approx \qty{22.5}{kg/mol}$.
\textbf{21.} Fraction en nombre $n_1 = 1 - p = 0.01$, soit 1~\% des molécules~; fraction en masse $w_1 = (1-p)^2 =
10^{-4}$, soit 0.01~\% de la masse.
\textbf{22.} Vers $x = -1/\ln 0.99 \approx 99.5$, c'est-à-dire près de $\bar X_n = 100$ (exercice 10).
\textbf{23.} $1000/226.32 = \qty{4.42}{mol}$ de \omterm{def:b2:polymer-synthesis:polymer}{motifs de répétition}, avec deux liaisons amide chacun~:
\qty{8.84}{mol} d'eau, soit $8.84 \times 18.015 \approx \qty{159}{g}$.
\textbf{24.} Des chaînes trop courtes donnent une fibre fragile~; des chaînes trop longues donnent une masse
fondue trop visqueuse pour être poussée à travers les fins orifices de la filière.
\textbf{25.} $\bar X_n = 20\,000/113.16 = 176.7$ et $p = 1 - 1/176.7 = \boldsymbol{\approx 0.994}$.
\end{solution}
